\chapter{三角圏}
\VerdierPhysicalPageAnchor{0110}%
\VerdierSemanticAnchor{ch2}
\begingroup
\small
\setlength{\abovedisplayskip}{4pt plus 2pt minus 2pt}%
\setlength{\belowdisplayskip}{4pt plus 2pt minus 2pt}%

\VerdierSemanticAnchor{ch2-1}
\section{定義と初歩的性質}

\VerdierSemanticAnchor{ch2-1-1}
\subsection{三角圏の定義}

\VerdierSemanticAnchor{ch2-1-1-1}
\noindent\textbf{定義 1.1.1.}\quad
\emph{三角圏} $\mathcal D$ とは、三角（第 I 章、3.2.1）の族を備えた厳密な
加法的 $\mathbb Z$-圏（第 I 章、1.5.11）であって、それらの三角を
\emph{完全三角}と呼び、次の性質を満たすものをいう。

\VerdierSemanticAnchor{ch2-1-1-1-tri}
\smallskip
\noindent\textbf{TRI:}
$\mathcal D$ の三角で完全三角に同型なものは、すべて完全三角である。任意の
対象 $X$（$\mathcal D$ の対象）に対し、三角
\[
  X\xrightarrow{\operatorname{id}_X}X\longrightarrow 0
  \longrightarrow X[1]
\]
は完全三角である。任意の射 $u\colon X\longrightarrow Y$（$\mathcal D$ の射）は、
次の完全三角に含まれる。
\[
  X\xrightarrow{u}Y\xrightarrow{v}Z\xrightarrow{w}X[1].
\]

\VerdierSemanticAnchor{ch2-1-1-1-trii}
\noindent\textbf{TRII:}
$\mathcal D$ の三角
\[
  X\xrightarrow{u}Y\xrightarrow{v}Z\xrightarrow{w}X[1]
\]
が完全三角であるための必要十分条件は、三角
\[
  Y\xrightarrow{v}Z\xrightarrow{w}X[1]
  \xrightarrow{-u[1]}Y[1]
\]
が完全三角であることである。

\VerdierSemanticAnchor{ch2-1-1-1-triii}
\noindent\textbf{TRIII:}
任意の二つの完全三角

\VerdierSemanticAnchor{ch2-1-1-1-triii-triangles}
\[
  \begin{aligned}
    X  &\xrightarrow{u}  Y  \xrightarrow{v}  Z  \xrightarrow{w}  X[1],\\
    X' &\xrightarrow{u'} Y' \xrightarrow{v'} Z' \xrightarrow{w'} X'[1]
  \end{aligned}
\]
と任意の可換図式

\VerdierSemanticAnchor{ch2-1-1-1-triii-square}
\begin{center}
  \begin{tikzcd}[column sep=4.5em,row sep=3em]
    X \arrow[r,"u"] \arrow[d,"f"'] & Y \arrow[d,"g"] \\
    X' \arrow[r,"u'"] & Y'\,,
  \end{tikzcd}
\end{center}
\endgroup
